\documentclass[11pt]{article}
\usepackage{amsmath,amssymb,amsthm,booktabs}
\usepackage[margin=2.4cm]{geometry}
\newtheorem{theorem}{Theorem}\newtheorem{lemma}[theorem]{Lemma}\newtheorem{proposition}[theorem]{Proposition}
\newtheorem{corollary}[theorem]{Corollary}
\theoremstyle{definition}\newtheorem{definition}[theorem]{Definition}\newtheorem{remark}[theorem]{Remark}
\newcommand{\Z}{\mathbb Z}
\title{Q597 G1c: the ascending neighbour of a swept tile is never swept\\
\normalsize PROVE 2026-11-20 c2. Step~0 (the $\psi$-separation test of the brief) \textbf{fails}, 0/1639. G1c is \textbf{proved} by another route:
$\psi$ alone gives the future half (Theorem~\ref{thm:F}, by hand from the local table), and a deterministic word model of journeys gives the whole statement (Theorem~\ref{thm:AW}, computer-assisted, exhaustive), with an independent hand proof for triangles (Proposition~\ref{prop:strip}).}
\author{Clio}\date{20 November 2026}
\begin{document}\maketitle

\section*{Provenance}
\begin{itemize}
\item Code: \texttt{proofs/code-1120c2-Q597/} (new this session; no script shadows a stdlib module). Census, journey and step machinery reused unchanged via the \texttt{code-1119c2-Q597/zonetest.py} preamble; journey dumps reused from \texttt{code-1120-Q597/dumpj5-*.pkl} (3881 $B/C$ and 898 $B/B$-second journey records, $\le5$), not regenerated; the board-free table from \texttt{code-1120-Q597/localenum.py}.
\item Scripts: \texttt{psisep.py} (Step~0), \texttt{threat.py}, \texttt{omtest.py} (threat sets, 4-offsets), \texttt{latstat.py}, \texttt{offsets.py}, \texttt{future.py}, \texttt{future2.py}, \texttt{fcontrol.py} (Theorem~\ref{thm:F} and its control), \texttt{lemmaD.py}, \texttt{lemmaD3.py}, \texttt{d1local.py}, \texttt{cross*.py} (the past, explored), \texttt{absmoves.py}, \texttt{words.py}, \texttt{words2.py} (Theorem~\ref{thm:AW}), \texttt{bbaw.py} ($B/B$ consistency), \texttt{fibre.py} (retreat). Outputs \texttt{*.out} beside them.
\item Upstream, all from my own files: \texttt{2026-11-20-Q597-no-repeat.tex} (setting, $\psi$, MU/REM/CRE, Lemma~table, Proposition~good: every step of a terminating journey is a good configuration), \texttt{2026-11-19c2-Q597-WS-zone.tex} (the A/B/T dictionary, Lemma~aneigh, gaps G1a--c). No external source is used in any proof. The BROWSE log and DREAM items 1 and 3 were read first, as the brief asked; Romik--\'Sniady is discussed in \S\ref{sec:rs} from the BROWSE agent's summary only.
\end{itemize}

\section*{Summary}
\begin{enumerate}
\item \textbf{Step 0 fails, 0/1639} (\S\ref{sec:step0}). For every (A)-tile $u$ of the 521 records, $\psi(u-\alpha)\le\mu_0$ componentwise: $\psi$ does \emph{not} separate the neighbour at the start of $J_o$. The preconditions hold 1639/1639, so the test failed for the right reason.
The reason is structural. $\psi$ sees only the forward direction: the rhombus sweeps from high $\psi$ to low $\psi$, and $u-\alpha$ is \emph{higher} than $u$ ($\psi(u-\alpha)=\psi(u)+e_c$). So in the past it sits exactly where a swept tile would sit. The threat set $\{k:\psi(u-\alpha)\le\mu_k\}$ is always an initial segment ending inside the run that removes $u$ (1639/1639).
\item \textbf{Theorem~\ref{thm:F} (the future half, by hand).} If a non-rhombus tile $u$ is removed at step $j$, then every tile $t$ with $\psi_c(t)\ge\psi_c(u)+1$ for some $c$ is removed at no step $k>j$. The new ingredient is Lemma~\ref{lem:W} (a dip of $\psi$ below $\mu$ forces $\mu$ down by one at that step). It is proved by hand from the three hexagon types.
\item \textbf{Theorem~\ref{thm:AW} (all of G1c).} If $u$ is removed at step $j$ and $v$ is an edge vector of $R_j$ with $\psi(u+v)=\psi(u)+e_c$ (an \emph{ascending} edge vector), then $u+v$ is removed at \emph{no} step of the journey.
The proof has three parts. (i) A journey is a word in $\{a0,a1,b\}$: each (class, type) has exactly one good configuration, so the word and $R_0$ determine every hexagon. (ii) A window lemma: the two removals would be at most 11 steps apart. (iii) An exhaustive check of all words of length $\le12$ (3\,188\,808 nodes, 33\,480\,864 step pairs, 0 violations). The control (ascending but \emph{not} an edge of $R_j$) fires 708\,600 times.
\item \textbf{Proposition~\ref{prop:strip}}: an independent hand proof for triangles. In the lattice of a $0$-triangle the swept cells form a path whose strip index increases by exactly one per cell and whose strip is weakly decreasing. The ascending translate of a path cell is never on the path.
\item \textbf{G1c follows} (Corollary~\ref{cor:G1c}). Its hypothesis, ``$-\alpha$ is an ascending edge vector of the rhombus at the step removing $u$'', is part of G1a: computed 1639/1639, not proved. Lemma~aneigh (11-19 c2) is no longer needed for G1c.
\item \textbf{Step 2}: G1a is scoped in \S\ref{sec:G1a}. \textbf{Retreat}: the forgetting fibre test is in \S\ref{sec:fibre}. \textbf{Romik--\'Sniady}: what $\psi$ gives and what is new is in \S\ref{sec:rs}; (z5) still needs running.
\end{enumerate}

\section{Setting}\label{sec:set}
As in \texttt{2026-11-20-Q597-no-repeat.tex} \S1. Points $P=(P_0,P_1)\in\Lambda_0\times\Lambda_1$, $\Lambda_i=\Z^2$, coordinates $(a,b)$. The target $T\in\{B,C\}$ has height $\eta_B(a,b)=b$, $\eta_C(a,b)=a+b$, and
\[\psi(t)=\Bigl(\min_{P\in t}-\eta_T(P_1),\ \min_{P\in t}\eta_T(P_0)\Bigr),\qquad \mu_k=\psi(R_k).\]
Proved there for every terminating journey in a Purbhoo mosaic: every step is a \emph{good} configuration of the table (Proposition~good). Hence MU ($\mu$ non-increasing), REM (removed non-rhombus tiles have $\psi\le\mu_k$), CRE, and no-revisit (Theorem~main) hold.
$R_k=s_0\times s_1$ with $s_0\subset\Lambda_0$, $s_1\subset\Lambda_1$ unit segments of one coordinate vector $\pm u_k$ (the \emph{class}), $\eta_T(u_k)=\pm1$ (Lemma~class with W1).
Translating a tile by $v\in\Lambda_0$ changes only $\psi_1$, by $\eta_T(v)$; by $v\in\Lambda_1$ only $\psi_0$, by $-\eta_T(v)$. So of the four edge vectors $\pm U_0,\pm U_1$ of $R_k$ exactly two are \emph{ascending} ($\psi(t+v)=\psi(t)+e_c$), one in each lattice.
Also put $\lambda_k=(\min_{P\in R_k}a(P_0),\ \min_{P\in R_k}a(P_1))$ and, for a tile, $A(t)$ likewise. The \emph{4-offset} is $\mathrm{om}(t,R)=(\psi(t)-\psi(R),\,A(t)-A(R))\in\Z^4$.

\section{Step 0: the brief's test fails}\label{sec:step0}
\texttt{psisep.py 5} (521 records, 1639 (A)-tiles, $q$: 825, $p$: 814).
\begin{itemize}
\item \textbf{Preconditions, asserted before counting:} $u-\alpha\in M_0$ in 1639/1639. At the step $j_r$ of $J_o$ that removes $u$, $-\alpha$ is an edge vector of $R_{j_r}$ in 1639/1639 ($q$: 572 type a, 253 type b; $p$: 546 / 268).
\item \textbf{The test} ``$\exists i:\ \psi_i(u-\alpha)>(\mu_k)_i$ for every $k$'' holds in \textbf{0/1639}. Already at $k=0$, $\psi(u-\alpha)\le\mu_0$ in every case. The defects $\psi(u-\alpha)-\mu_0$ are $(0,-1),\dots,(0,-5)$ on side $q$ and $(0,0),(-1,0)$ on side $p$.
\item \textbf{Lesion} ($\eta_B\leftrightarrow\eta_C$): also 0/1639. Since the test itself fails, the lesion is uninformative, and I say so.
\item \textbf{Control ($B/B$).} On $B/B$ ``A-shaped'' tiles, $\psi$-separation at $k=0$ holds 758 times, and then $u-\alpha$ is never removed (758/758, as REM predicts). It fails 1863 times, and in 1097 of those $u-\alpha$ \emph{is} removed. The brief's prediction ``the test should still pass on $B/B$'' was wrong because the test does not pass on $B/C$ either.
\item \textbf{The prefix:} $\psi$ of $y$'s prefix journey does not separate either (0/1639). The prefix never removes $u-\alpha$ (as on 11-19 c2).
\end{itemize}
\paragraph{What the failure says} (\texttt{threat.py 5}). Always $\psi(u-\alpha)=\psi(u)+e_c$, where $c=1$ on side $q$ ($-\alpha\in\Lambda_0$; tiles: squares and $0$-triangles only) and $c=0$ on side $p$ ($-\alpha\in\Lambda_1$; squares and $1$-triangles only). So $-\alpha$ is an \emph{ascending} edge vector of $R_{j_r}$, 1639/1639.
The threat set $K=\{k:\psi(u-\alpha)\le\mu_k\}$ is an initial segment $[0,k^*]$ with $k^*\le$ the last step of $u$'s run, 1639/1639. So $\psi$ handles every step after the run, and nothing before it.
That is forced. REM says the rhombus sweeps downward in $\psi$, and $u-\alpha$ lies above $u$. A tile above $u$ is exactly what the rhombus \emph{would} have swept earlier. The past needs lateral information that $\psi$ does not carry.
The 4-offset does not suffice either (\texttt{omtest.py}). $\mathrm{om}(u-\alpha,R_k)$ lies in the removal set $Q_T$ (below) at some $k$ for 1424 tiles: within the run (type a), or at $k=j_r-1$ (type b, 383 tiles). Yet $u-\alpha$ is never removed.

\section{The future half, by hand}\label{sec:F}
\begin{lemma}[dip]\label{lem:W}
At a good step $k$, every tile $t$ of $H_k$ has $\psi_c(t)\ge(\mu_{k+1})_c$ and $\psi_c(t)\ge(\mu_k)_c-1$ for $c=0,1$. If $\psi_c(t)<(\mu_k)_c$ for some tile $t$ of $H_k$, then $\psi_c(t)=(\mu_k)_c-1=(\mu_{k+1})_c$.
\end{lemma}
\begin{proof}
Take $c=1$; then $\psi_1(t)=\min\eta(\pi_0 t)\ge m:=\min\eta(\pi_0H_k)$, since $\pi_0t\subseteq\pi_0H_k$. The three hexagon types of Lemma~hex (used in Lemma~table) are as follows.
\emph{a1}: $H=Z(U_0,U_1,G)$ with $G\in\Lambda_1$, so $\pi_0H=s_0$ and $m=(\mu_k)_1$. There is no dip, and $\pi_0R_{k+1}=s_0$.
\emph{b}: $\pi_0H=\tau$, a unit triangle with $s_0$ as an edge. A unit triangle has one edge in each of the three classes, so exactly one edge lies in $\ker\eta$. Its $\eta$-values are therefore the two values $(\mu_k)_1,(\mu_k)_1+1$ taken on $s_0$, so $m=(\mu_k)_1$. Again no dip, and $\pi_0R_{k+1}\subset\tau$ has $\min\eta=(\mu_k)_1$ as well.
\emph{a0}: $\pi_0H=Z(U_0,G)$, a unit parallelogram, $G\in\Lambda_0$. The step is the point reflection of $H$ about its centre $z$, so $\pi_0R_{k+1}=2z_0-s_0$, the side of $\pi_0H$ opposite to $s_0$, which is $s_0\pm G$.
The minimum of the linear function $\eta$ over the parallelogram is attained on one of its two $U_0$-sides, so $m=\min((\mu_k)_1,(\mu_{k+1})_1)$, and $m\ge(\mu_k)_1-|\eta(G)|\ge(\mu_k)_1-1$.
In all cases $m=\min((\mu_k)_1,(\mu_{k+1})_1)\ge(\mu_k)_1-1$, and $(\mu_{k+1})_1\le(\mu_k)_1$ by MU. If $\psi_1(t)<(\mu_k)_1$, then $m<(\mu_k)_1$, so $(\mu_{k+1})_1=m$, and $(\mu_k)_1-1\le m\le\psi_1(t)\le(\mu_k)_1-1$.
For $c=0$ exchange the lattices (a0$\leftrightarrow$a1).
\end{proof}
\noindent\emph{Check} (\texttt{future.py}): on the dumps, every dip has $\psi_c=\mu_c-1$ and $\mu$ drops at that step, 28390/28390 ($B/C$) and 7696/7696 ($B/B$-second). On the table, the $\Delta\mu$ column of Lemma~table agrees.

\begin{theorem}[removal is $\psi$-antitone]\label{thm:F}
In a terminating journey in a Purbhoo mosaic, let $u$ be removed at step $j$ and $t$ at step $k>j$ (both non-rhombus). Then $\psi(t)\le\mu_k\le\mu_{j+1}\le\psi(u)$ componentwise. In particular, a tile $t$ with $\psi_c(t)\ge\psi_c(u)+1$ for some $c$, for instance the ascending translate $u+v$ by any vector $v$, is removed at no step after $j$.
\end{theorem}
\begin{proof}
REM gives $\psi(t)\le\mu_k$. MU gives $\mu_k\le\mu_{j+1}$. Lemma~\ref{lem:W} at step $j$ gives $\psi_c(u)\ge(\mu_{j+1})_c$ for both $c$.
\end{proof}
\noindent\emph{Check and control} (\texttt{future2.py}, \texttt{fcontrol.py}). Over all dumped journeys, all steps $j$, all removed $u$ and all 12 unit vectors $w$:
\begin{itemize}
\item ascending translates are removed later \textbf{0} times (356\,832 $B/C$ + 85\,284 $B/B$-second);
\item flat translates are removed later 11\,902 + 1\,784 times, and descending translates 17\,179 + 4\,589 times.
\end{itemize}
The detector is live. The statement depends only on ascent, not on $w$ being an edge of $R_j$, as the theorem says.

\section{The whole statement: the word model}\label{sec:AW}
\begin{lemma}[a journey is a word]\label{lem:word}
For each $T$ and each of the two allowed classes, the table has exactly one good configuration of each type a0, a1, b. These are sign-consistent: classes $u$ and $-u$ give translates.
Consequently, in a terminating journey the step types $w_k\in\{a0,a1,b\}$, together with $R_0$, determine every $R_k$, every hexagon $H_k$, and the tiles removed and created at step $k$.
\end{lemma}
\begin{proof}
Read off Lemma~table (no-repeat tex): for $B$, class $(1,-1)$ has a0, a1, b each once (across edge undefined), and class $(0,1)$ has a0 interior, a1 interior, b interior once each, with the boundary ones bad. $C$ is the same. \texttt{words.py} asserts uniqueness and sign-consistency on the enumerated table.
By Proposition~good every step is good, so step $k$ is the unique good configuration of (class of $R_k$, $w_k$), translated to $R_k$. Induct on $k$.
\end{proof}
The abstract moves (\texttt{absmoves.out}), with $\Delta(\mu,\lambda)$ per configuration (\texttt{words4.out}):
\begin{center}\begin{tabular}{lll}\toprule
$T$ & type & $\Delta(\mu_0,\mu_1,\lambda_0,\lambda_1)$ for the two classes\\\midrule
$B$ & a0 & $(0,-1,1,0)$, $(0,0,1,0)$\\
 & a1 & $(-1,0,0,0)$, $(0,0,0,1)$\\
 & b & $(0,0,0,0)$ [from class $(0,1)$], $(0,0,1,1)$ [from $(1,-1)$]\\
$C$ & a0 & $(0,-1,0,0)$, $(0,0,1,0)$\\
 & a1 & $(-1,0,0,1)$, $(0,0,0,1)$\\
 & b & $(0,0,0,0)$, $(0,0,1,1)$\\\bottomrule
\end{tabular}\end{center}
So $\mu$ is non-increasing and $\lambda$ non-decreasing (computed on the dumps too, \texttt{latstat.out}: $\Delta\min a\in\{0,1\}$ at every step of 3881 + 898 journeys).
Only a b-step can be \emph{stationary} ($\Delta(\mu,\lambda)=0$). A b-step changes the class, and the b-step from the other class is not stationary, so \textbf{no two consecutive steps are stationary}.
The removed tiles have 4-offsets in
\[Q_B=\{(0,0,0,0),(-1,0,0,0),(0,-1,0,0),(0,0,1,0),(-1,0,1,0)\},\quad Q_C=\{(0,0,0,0),(-1,0,0,0),(0,-1,0,0),(0,0,0,1),(0,-1,0,1)\}\]
(table, \texttt{words4.out}; the same sets occur on all dumped journeys, \texttt{offsets.out}).

\begin{lemma}[window]\label{lem:window}
Let $u$ be removed at step $j$ and $t=u+v$ at step $k\neq j$, with $v$ an ascending edge vector of $R_j$. Then $|j-k|\le11$.
\end{lemma}
\begin{proof}
Put $f(i)=\mathrm{om}(u,R_i)$. Each coordinate of $f$ is monotone in $i$: the $\psi$-part is non-decreasing because $\mu$ is non-increasing, and the $A$-part is non-increasing. So the total variation of $f$ between $k$ and $j$ equals $|f(j)-f(k)|_1$.
Every non-stationary step changes $f$ by at least $1$ in $\ell^1$. We have $f(j)\in Q_T$, and $f(k)=\mathrm{om}(t,R_k)-\delta\in Q_T-\delta$, where $\delta=\mathrm{om}(t,R)-\mathrm{om}(u,R)$ is the fixed shift of $v$. For ascending edge vectors, $\delta\in\{(0,1,0,0),(0,1,-1,0),(1,0,0,0),(1,0,0,1)\}$ for $B$ and $\{(0,1,1,0),(0,1,0,0),(1,0,0,-1),(1,0,0,0)\}$ for $C$.
The maximum of $|q-(q'-\delta)|_1$ over $q,q'\in Q_T$ and these $\delta$ is $5$ for both $T$. So at most $5$ of the $|j-k|$ steps between $k$ and $j$ are non-stationary. Stationary steps are never consecutive, so $|j-k|\le 2\cdot5+1=11$.
\end{proof}

\begin{theorem}[the ascending neighbour of a swept tile is never swept]\label{thm:AW}
In a terminating journey toward $B$ or $C$ in a Purbhoo mosaic, let $u$ be a non-rhombus tile removed at step $j$, and let $v$ be an ascending edge vector of $R_j$. Then $u+v$ is removed at no step $k\neq j$. Also $u+v\notin H_j$.
\end{theorem}
\begin{proof}
By Lemma~\ref{lem:window}, $k$ and $j$ lie in a window of at most $12$ consecutive steps. By Lemma~\ref{lem:word}, the removals in that window are those of a word of length $\le12$ started from $R_{\min(j,k)}$, which has one of the two allowed classes. Everything is invariant under translation.
\texttt{words2.py 12} enumerates every word of length $\le12$ from each allowed class, for both $T$. For every pair of steps it tests whether an ascending $R$-edge translate of a tile removed at one step is removed at the other: \textbf{0 violations} in 33\,480\,864 pairs (3\,188\,808 nodes).
\emph{Second, independent implementation} (\texttt{words3.py 12}). Its 12 configurations come from the \emph{engine}: they are extracted from the 36\,843 real steps and asserted identical across all occurrences. Ascent is recomputed from $\eta_T$, and the search is an iterative DFS with its own data structures. It finds \textbf{0 violations} in 33\,480\,792 pairs (3\,188\,640 nodes).
The counts differ slightly from \texttt{words2}, which re-counts the 3-letter job prefixes: $168$ nodes and $72$ pairs.
The window ingredients ($Q_T$, the $\Delta(\mu,\lambda)$ table, and which b-step is stationary) were re-derived from the engine configurations too (\texttt{windowcheck.out}), and they agree.
The step $k=j$ is covered by the table directly: in all 12 configurations the ascending translates of the removed tiles are disjoint from the removed tiles (\texttt{words2.prep}).
\end{proof}
\noindent\textbf{Controls.}
\begin{itemize}
\item The same enumeration with $v$ ascending but \emph{not} an edge vector of $R_j$ finds coincidences 708\,600 times. The test can fire, and the edge hypothesis is what matters for the \emph{past}. Theorem~\ref{thm:F} has already removed it for the future.
\item On real boards (\texttt{future2.py}), ascending $R_j$-edge translates of removed tiles are removed \emph{earlier} 0 times and later 0 times, over 221\,058 pairs. This is independent corroboration from the engine rather than the table.
\item \texttt{cross*.py}: of the past pairs where neither $\psi$ nor the 4-offset excludes the removal, 1884 are same-class and 3798 cross-class (an odd number of b-steps between). All of them lie in windows of the shapes aa, ab, aba, abaa, abaaa, abab, abaab, abb. This is consistent with Lemma~\ref{lem:window}.
\end{itemize}

\begin{proposition}[triangles, by hand]\label{prop:strip}
For $T=B$ and $0$-triangles (and so, by the reflection $\Sigma$ of the no-repeat tex, for $T=C$ and $1$-triangles), Theorem~\ref{thm:AW} has the following hand proof.
\end{proposition}
\begin{proof}
A $0$-triangle $\tau\times\{q\}$ lies in $H_k$ only if its cell $\tau$ lies in $\pi_0H_k$. a1-steps contain no $0$-triangles, and a0- and b-steps contain exactly the $0$-cells listed below (\texttt{absmoves.out}). So it suffices to show that the cell $\tau+v$ is never swept.
Write $\Delta(q)=\{q,q+(1,0),q+(0,1)\}$ and $\nabla(q)=\{q+(1,0),q+(0,1),q+(1,1)\}$. Both lie in the strip $q_b\le b\le q_b+1$, the \emph{strip} $q_b$. Index them by $\iota\Delta(q)=2q_a$ and $\iota\nabla(q)=2q_a+1$. Cells of one strip with consecutive indices share an edge.
The allowed classes for $B$ are $V=(0,1)$ and $D=(1,-1)$. Write $V@p=[p,p+(0,1)]$, which lies between $\iota=2p_a-1$ and $2p_a$ in strip $p_b$, and $D@p=[p,p+(1,-1)]$, which lies between $2p_a$ and $2p_a+1$ in strip $p_b-1$. The table gives:
\begin{itemize}
\item $V@p$: a0 sweeps $\Delta(p),\nabla(p-(0,1))$ and goes to $V@(p+(1,-1))$; b sweeps $\Delta(p)$ and goes to $D@(p+(0,1))$.
\item $D@p$: a0 sweeps $\nabla(p-(0,1)),\Delta(p+(1,-1))$ and goes to $D@(p+(1,0))$; b sweeps $\nabla(p-(0,1))$ and goes to $V@(p+(1,-1))$.
\end{itemize}
Read in (strip, index), every move sweeps the cells with the next one or two indices. The strip drops by one only inside $V$-a0, between its two cells: $\Delta(p)$ at $(p_b,2p_a)$, then $\nabla(p-(0,1))$ at $(p_b-1,2p_a+1)$.
The journey starts at $s_0$ of class $V$ (Lemma~class). Hence the swept cells form a sequence whose index increases by exactly $1$ per cell and whose strip is weakly decreasing. So the strip is a function $S(\iota)$ of the index on an interval of indices, and it is weakly decreasing. A strip is entered either at the start, at an even index, or by a drop, at an odd index.
The ascending vector is $v=(0,1)$ for class $V$ and $v=(-1,1)$ for class $D$. Translation by $v$ raises the strip by one and changes $\iota$ by $0$, resp.\ $-2$.
\emph{Class $V$ at step $j$} (cells $(s,2p_a)$, and for a0 also $(s-1,2p_a+1)$, with $s=p_b$): the translates are $(s+1,2p_a)$ and $(s,2p_a+1)$, but $S(2p_a)=s$ and $S(2p_a+1)=s-1$.
\emph{Class $D$ at step $j$} ($s=p_b-1$; cells $(s,2p_a+1)$, and for a0 also $(s,2p_a+2)$): the translates are $(s+1,2p_a-1)$ and $(s+1,2p_a)$. $D@p$ is not the start (the start is class $V$), so the previous move swept $\Delta(p-(0,1))$ at $(s,2p_a)$: $S(2p_a)=s$.
Strip $s$ was entered at an index $\le2p_a$. If at $2p_a$ (the start), index $2p_a-1$ is not swept at all. Otherwise it was entered at an odd index $\le2p_a-1$, so $S(2p_a-1)=s$.
In every case the translate is not a swept cell. Since this is a statement about the whole sequence, it covers all $k$, past and future.
\end{proof}
\noindent The same picture reproves no-repeat for $0$-triangles: indices increase strictly.
$\Sigma$ exchanges $(B,\Lambda_0)\leftrightarrow(C,\Lambda_1)$ and $(B,\Lambda_1)\leftrightarrow(C,\Lambda_0)$. So the proposition covers $0$-triangles for $B$ and $1$-triangles for $C$.
Those are exactly the triangle (A)-tiles that occur: side $q$ ($J_o\to B$) has only $0$-triangles and squares, and side $p$ ($J_o\to C$) only $1$-triangles and squares (\texttt{threat5.out}, computed).
\begin{proposition}[the mirrored case, by hand]\label{prop:strip2}
The conclusion of Proposition~\ref{prop:strip} also holds for $T=B$ and $1$-triangles (and so, by $\Sigma$, for $T=C$ and $0$-triangles). Hence Theorem~\ref{thm:AW} holds for every triangle tile by hand.
\end{proposition}
\begin{proof}
Now the walk is that of $s_1$ in $\Lambda_1$. The ascending vectors are those with $-\eta_B(v)=+1$: $v=(0,-1)$ for class $V$ and $v=(1,-1)$ for class $D$. The table (\texttt{absmoves.out}) gives:
\begin{itemize}
\item $V@p$: a1 sweeps $\Delta(p),\nabla(p)$ and goes to $V@(p+(1,0))$; b sweeps $\Delta(p)$ and goes to $D@(p+(0,1))$.
\item $D@p$: a1 sweeps $\nabla(p-(0,1))$, then $\Delta(p)$, and goes to $D@(p+(0,1))$; b sweeps $\nabla(p-(0,1))$ and goes to $V@(p+(1,-1))$.
\end{itemize}
Use $\kappa=\iota+2\sigma$, where $\sigma$ is the strip. Within a strip $\kappa$ rises by $1$ per cell. The one strip change is the \emph{climb} inside $D$-a1, from $\nabla(p-(0,1))$ at $(\sigma,2p_a+1)$ to $\Delta(p)$ at $(\sigma+1,2p_a)$. There $\kappa$ also rises by $1$.
So the swept cells have $\kappa$ increasing by exactly $1$ per cell, and the strip $S(\kappa)$ is weakly \emph{increasing}. A climb always lands on a $\Delta$.
Translation by $v$ lowers the strip by one. It changes $\kappa$ by $-2$ for class $V$ ($\iota$ unchanged) and by $0$ for class $D$ ($\iota+2$).
\emph{Class $D$ at step $j$}: the translate of a cell at $(\sigma,\kappa)$ is $(\sigma-1,\kappa)$, and $S(\kappa)=\sigma$.
\emph{Class $V$ at step $j$}, $s_1=V@p$, $\sigma=p_b$, $\kappa_0=\kappa(\Delta(p))$. The swept cells are $\Delta(p)$ at $\kappa_0$, and for a1 also $\nabla(p)$ at $\kappa_0+1$. Their translates sit at $(\sigma-1,\kappa_0-2)$ and $(\sigma-1,\kappa_0-1)$.
If $V@p$ is the start, no cell has $\kappa<\kappa_0$. Otherwise $V@p$ was reached by $V$-a1 from $V@(p-(1,0))$, which swept $\Delta(p-(1,0))$ and $\nabla(p-(1,0))$ in strip $\sigma$ at $\kappa_0-2$ and $\kappa_0-1$. Or it was reached by $D$-b from $D@(p-(1,-1))$, which swept $\nabla(p-(1,0))$ at $(\sigma,\kappa_0-1)$. That $D$ is not the start (the start has class $V$), and every move ending at a $D$ sweeps last the $\Delta$ just before it, here $\Delta(p-(1,0))$ at $(\sigma,\kappa_0-2)$.
Either way $S(\kappa_0-2)=S(\kappa_0-1)=\sigma\ne\sigma-1$, so no translate is swept.
\end{proof}
\noindent Only squares rest on the enumeration alone. The two hand proofs are out-of-sample checks of \texttt{words2}/\texttt{words3} on the triangle part.

\begin{corollary}[G1c]\label{cor:G1c}
Fix a clean nontrivial $B/C$ record. Let $u$ be an (A)-tile, removed by step $j_r$ of $J_o$. Assume \textup{(G1a-A)}: $-\alpha$ is an ascending edge vector of $R_{j_r}$. Then no step of $J_o$ removes $u-\alpha$. If moreover $u-\alpha\in M_0$, then $u-\alpha\in M_0\cdot J_o$.
\end{corollary}
\begin{proof}
Theorem~\ref{thm:AW} with $v=-\alpha$.
\end{proof}
\noindent(G1a-A) and $u-\alpha\in M_0$ hold in 1639/1639 (\S\ref{sec:step0}). They are part of G1a, not proved. Lemma~aneigh is superseded for this purpose. It used the run structure, and the corollary covers all runs and turns at once.
\emph{Not covered:} the common prefix of $y$, which is a different journey on $M_0\cdot J_o$. Prop.~wsb (A$^*$) needs it, and it is computed only (0 removals, \S\ref{sec:step0}).

\paragraph{Consistency on $B/B$}\label{sec:bbaw} (\texttt{bbaw.py 5}). The theorem predicts that on $B/B$, where WS fails, an A-shaped tile whose $u-\alpha$ \emph{is} removed by $J_o$ must violate (G1a-A). This is a prediction that could have failed.
Of the 2621 $B/B$ A-shaped tiles with $u-\alpha\in M_0$, all \textbf{1097} whose $u-\alpha$ is removed violate (G1a-A), and all \textbf{1524} that satisfy (G1a-A) keep $u-\alpha$. There are no other cases.
So on $B/B$ the failure of survival is exactly a failure of the two-journey dictionary (G1a), never of the single-journey theorem. This matches 11-19 c2's ``$B/B$ failure, located''.

\paragraph{Size 6 (fixed subsample, not the census)} (\texttt{aw6.py 6 3 24 3}: chunks 0--2 of 24, the same subsample as \texttt{pot6} of the no-repeat tex; 3951 states, 2681 distinct $B/C$ journeys, 794 $B/B$-second journeys).
Lemma~\ref{lem:word}: 30\,488/30\,488 steps match exactly one configuration. Theorem~\ref{thm:AW}: 0 violations in 182\,928 pairs. Theorem~\ref{thm:F}: ascending translates are removed later 0 times out of 365\,856. The control fires: flat 12\,876, descending 19\,367.

\paragraph{Lemma~\ref{lem:word} on real journeys} (\texttt{wordreal.py}). Every one of the 36\,843 steps of the dumped journeys (3881 $B/C$, 898 $B/B$-second) matches \emph{exactly one} abstract configuration of its (class, engine type), translated to $R_k$. This covers both the removed set and the new rhombus.

\section{Step 2: what G1a needs}\label{sec:G1a}
With G1b (no-revisit) and G1c (Corollary~\ref{cor:G1c}) proved, WS-b reduces to statements about \emph{two} journeys. These are what remains:
\begin{description}
\item[(G1a-0) Prefix disjointness.] $y$'s common prefix touches neither $u$ nor $u-\alpha$. Given this, $u\in M_P\setminus Mx_P$ forces $u\in M_0$ removed by $J_o$. This is pure bookkeeping, and it is the part that turns ``read wake'' into ``$x$'s wake''.
\item[(G1a-1) The $\alpha$-dictionary.] If $u$ is read by $y$'s crossing and removed by step $j_r$ of $x$, then $-\alpha$ (a displacement of \emph{$y$'s} rhombus) is an edge vector of $x$'s rhombus $R_{j_r}$. It is ascending when $u$ is of kind (A), and it equals the displacement of $u$ under step $j_r$ when $u$ is of kind (B).
This links the two journeys. It is the real content of G1a, and of the Romik--\'Sniady-type statement (z5).
\item[(G1a-2) The trichotomy.] Every read wake tile is A, B or T. Given (G1a-1), this is a statement about which edge class of $u$ is parallel to $\alpha$, and it is local to step $j_r$.
\item[(G1a-3)] For (A): $u-\alpha\in M_0$. For (B)/(T): $u-\alpha$ is made by $x$. Then G1b and G1c give survival through $J_o$.
\end{description}
\paragraph{(G1a-A) splits in two (computed, \texttt{g1aparity.py 5}, 521 records).}
\begin{itemize}
\item[(A1)] $-\alpha$ is an edge vector of $x$'s \emph{starting} rhombus $R_0$: 1639/1639.
\item[(A2)] Every (A)-tile is removed at a step $j_r$ preceded by an \emph{even} number of b-steps of $J_o$, i.e.\ while $R$ still has its starting class: 1639/1639 ($q$: 572 a + 253 b; $p$: 546 a + 268 b).
\end{itemize}
Since the class is constant between b-steps and ascent depends only on $v$, (A1) and (A2) together imply (G1a-A).
Moreover (A1) is a statement about the start only. With the nest-$A$ frame $u_A=(0,-1;0,0)$, $w_A=(0,0;0,1)$ (no-repeat tex, Lemma~start), the ascending edge vectors of $R_0$ are $-u_A$ and $-w_A$ for both $B$ and $C$ ($\eta_B(0,1)=\eta_C(0,1)=1$). So (A1) together with the computed ascent (1639/1639) implies \emph{$\alpha\in\{u_A,w_A\}$, a frame vector of nest $A$}. This looks like the most tractable piece of G1a, since it concerns $y$'s displacement and the frame only. (A2) is the genuinely two-journey part.
For the (B)/(T) tiles the parity is even in 2654 cases and odd in 51 (all type-b removals). So the parity statement is specific to (A).

\begin{corollary}[WS-b reduces to G1a; bookkeeping]\label{cor:wsb}
Assume, for a read wake tile $u$ of a clean nontrivial $B/C$ record, that (G1a-0) holds and that $u$ is of kind (A) with (G1a-1) and (G1a-3), or of kind (B)/(T) with (G1a-1) and (G1a-3). Then WS-b holds at $u$.
\end{corollary}
\begin{proof}
(A): Corollary~\ref{cor:G1c} gives $u-\alpha\in M_0\cdot J_o$, and (G1a-0) keeps it through the prefix. So (A$^*$) of Prop.~wsb holds.
(B)/(T): $u-\alpha$ is made by a step of $J_o$ and is not removed later by Theorem~main (no-revisit) of the no-repeat tex. (G1a-0) covers the prefix. So (B$^*$) holds.
Prop.~wsb (11-19 c2) concludes.
\end{proof}
So the next proof target is (G1a-1), a statement relating the rhombi of $x$ and $y$ at a read tile. $\psi$ and the word model are single-journey tools, and they cannot give it alone.

\section{Retreat: the forgetting fibre test}\label{sec:fibre}
From \texttt{memory/connections/2026-11-20-a-common-third-step-is-a-forgetting-step.md}, falsifiers 1 and 2. \texttt{fibre.py 5} covers all 440 non-commuting $B/B$ pairs at $\le5$: the 162 \emph{nonvacuous} pairs (nest $A$ keeps $\ge2$ rhombi after $x,y$) and the 278 others. Put $D=xy(s)\triangle yx(s)$, with $c_{12}=xy(s)$ and $c_{21}=yx(s)$.
\begin{center}\begin{tabular}{lrr}\toprule
statement & nonvacuous & vacuous\\\midrule
(a) $D\subseteq\mathrm{wake}(x)\cup\mathrm{wake}(y)$ (first-journey wakes $s\triangle x(s)$, $s\triangle y(s)$) & \textbf{0}/162 & 0/278\\
(a$'$) $D\subseteq$ the four journey wakes & 162/162 & 278/278\\
(b) union form, $D\subseteq(c_{12}\triangle c_{12}m)\cup(c_{21}\triangle c_{21}m)$ & 171/171 & 278/278\\
(b) per side, $D\subseteq c_{12}\triangle c_{12}m$, or $D\subseteq c_{21}\triangle c_{21}m$ & \textbf{0}/171, 0/171 & 0/278\\
closing label $m$ unique / two labels & 153 / 9 & 278 / 0\\
target of every closing $m$ & $B$ & $B$\\\bottomrule
\end{tabular}\end{center}
\textbf{Instrument honesty.} (a$'$) and the union form of (b) are \emph{tautologies}. Indeed $c_{12}\triangle c_{21}\subseteq(s\triangle c_{12})\cup(s\triangle c_{21})$, and if $c_{12}m=c_{21}m=F$, every $t\in D$ is changed by $m$ on exactly one side. The second fact also makes $|D\cap m_{12}|+|D\cap m_{21}|=|D|$ automatic, and it is seen in every row of \texttt{fibre5.out}. They passed as instrument checks, not as evidence.
\begin{itemize}
\item \textbf{Falsifier 1, verdict: the prediction fails in both of its contentful forms.} $D$ always leaves the union of the first-journey wakes: the order-dependence is carried by the \emph{second} journeys. And neither side's $m$ touches all of $D$: $m$ removes part of $D$ on one side and re-creates the rest on the other, splitting $|D|=32,\dots,80$ unevenly (e.g.\ $13+19$, $25+43$).
So ``$m$ forgets exactly the order'' is true only in the tautological union sense. The non-injectivity of $m$ (the connection's actual point) stands. What is refuted is the stronger reading that $m$'s wake on one board covers the difference.
\item \textbf{Falsifier 2 (Hecke/braid $x\cdot y\cdot x=y\cdot x\cdot y$): ill-posed as a label identity.} A migration label is a nest-$A$ corner, and $x$'s own migration vacates $x$. So ``$x$ again'' after $x,y$ is undefined except by accident, and \texttt{fibre.py} found it defined and equal in 0/440 cases.
The 11-20 table already gives the honest form: closures with \emph{different} labels on the two sides occur 0 times out of 440. The braid shape does not occur in any form.
\item \textbf{New fact:} the closing migration is always toward $B$ (171/171 nonvacuous and 278/278 vacuous), never toward $C$. In 9 pairs two different $B$-labels close.
\end{itemize}

\section{Romik--\'Sniady Prop.~3.1: what $\psi$ already gives}\label{sec:rs}
From the BROWSE summary only (agent-summary level; the paper was not read this session). Prop.~3.1 has three parts: a slide changes edges only in the tree of its path; the changes are at W/S neighbours of the path; each change is in a fixed direction.
\begin{itemize}
\item \emph{Fixed direction}: given by the direction law (Corollary~dl of the no-repeat tex). Each moved triangle moves $+1$ in its own height.
\item \emph{One side of the route}: given by REM/CRE with MU. Removed tiles lie at or below $\mu_k$, created ones strictly above $\mu_{k+1}$ in some coordinate. Theorem~\ref{thm:F} sharpens this: removal is $\psi$-antitone.
\item \emph{New today}: the rim is untouched. Theorem~\ref{thm:AW} says that the ascending neighbour of anything swept is never swept, at any time. This is the mosaic form of ``changes are adjacent to the path, and the path's own rim survives''. Proposition~\ref{prop:strip} draws it as a strip-monotone path.
\item \emph{Not given}: (z5) compares the wake of \emph{one} slide with the route of \emph{another} slide, after translation by $\alpha$. That is exactly (G1a-1), a two-journey statement, so \textbf{(z5) still needs running.}
\end{itemize}

\section{Where the proof stops}\label{sec:gaps}
\begin{description}
\item[Computer-assisted parts.] Lemma~\ref{lem:word} (uniqueness per (class, type), read off the table), the sets $Q_T$ and the $\Delta(\mu,\lambda)$ table, and the word enumeration of Theorem~\ref{thm:AW}. All are exhaustive and exact on integer points. They rest on the same board-free table as no-repeat (two enumerators there).
\emph{Enumerators.} The hexagon sets behind the 12 configurations are confirmed by the independent BFS \texttt{localenum2.py} (36/36, no-repeat tex). The word enumeration was done twice: \texttt{words2.py} on \texttt{localenum.analyse}'s configurations, and \texttt{words3.py} on configurations extracted from the engine. Both found 0 violations. Further checks against the engine: Lemma~\ref{lem:word} matches 36\,843/36\,843 real steps, and the conclusion of Theorem~\ref{thm:AW} holds on 221\,058/221\,058 real pairs.
\item[By hand.] Lemma~\ref{lem:W}, Theorem~\ref{thm:F}, Lemma~\ref{lem:window} (given the tables), and Propositions~\ref{prop:strip} and~\ref{prop:strip2} (Theorem~\ref{thm:AW} for all triangles). Squares are by enumeration only.
\item[Hypotheses.] Termination, Purbhoo mosaic (Lemma~N), and the engine's rule (S1)--(S4), as for no-repeat. G1c itself needs (G1a-A), which is computed.
\item[Not proved.] G1a, the prefix part of (A$^*$), WS, CT, BS.
\end{description}

\section*{Status}
\begin{itemize}
\item \textbf{Proved:} Lemma~\ref{lem:W}, Theorem~\ref{thm:F} (by hand), Theorem~\ref{thm:AW} (computer-assisted), Propositions~\ref{prop:strip}, \ref{prop:strip2} (by hand), and Corollary~\ref{cor:G1c} (G1c given (G1a-A)).
\item \textbf{Refuted:} the brief's Step~0 mechanism, ``$\psi$ separates $u-\alpha$ from all of $J_o$'' (0/1639).
\item \textbf{Computed:} (G1a-A) 1639/1639; threat sets; the controls above; the fibre test (\S\ref{sec:fibre}).
\end{itemize}
\end{document}
